\section{Photographs of the Final Controller}\label{app:photos}
\Fref{fig:board} documents the soldered Veroboard implementation described in \Sref{sec:hardware}. The board has three eight-pin op-amp packages and a blue bias trimmer. The photograph records the final physical assembly; it does not establish every underside connection or measured resistance.
\fig{../../hardware/veroboard.jpeg}{.84}{Supplied final Veroboard photograph. The three op-amp stages and adjustable bias trimmer are visible.}{fig:board}
The circuit design specifies 741-type devices and bipolar supply rails. Device markings and decoupling connections should be checked against the physical board when verifying the as-built netlist. No photographs of handwritten calculations are used in this report.
\page

\section{Experimental Demonstration Record}\label{app:demo}
The original display in \Fref{fig:demo} supports the response evaluated in \Sref{sec:tests}. The red trace is the reference and the blue trace is altitude. This image is also retained in \texttt{data/session2/hardware\_step\_test.jpeg}.
\fig{../../data/session2/hardware_step_test.jpeg}{1}{Original controller-test display for PTRJUN003, showing the altitude response and terminal readouts.}{fig:demo}
The numerical values transcribed in \Tref{tab:display} support the tracking and sensor checks.
\begin{table}[H]\centering\small
\caption{Numerical readouts transcribed from the retained test display.}\label{tab:display}
\begin{tabular}{lr}\toprule\thead
Field & Displayed value\\\midrule
Student number & PTRJUN003\\
Set point & 5.033 m\\
Height & 4.897 m\\
Tracking error & 0.137 m\\
Angle of attack / command & 2.498 V\\
Height (left panel) & 4.896 m\\
Sensor output (left panel) & 4.553 V\\
Elapsed time & 86.10 s\\\bottomrule
\end{tabular}\end{table}
The 0.001 m differences between the readouts are consistent with rounding or asynchronous display updates. The assessor's marking sheet and signature were not supplied, so an awarded demonstration mark is not claimed.
\page

\section{Original Simulink Evidence}\label{app:simulink}
\Fref{fig:originalmodel} preserves the supplied model image. Its visible values match the identified plant, sensor coefficient and selected gain used in \Sref{sec:design}.
\fig{../../data/session2/simulink_model.png}{1}{Supplied Simulink model: $k_s=0.93$, $K_p=0.01432$, $u_h=2.5$ V, velocity plant $14.2748/(8.1s+1)$ and position integrator.}{fig:originalmodel}
\Fref{fig:originalscope} begins at a 2.5 m reference while the output starts near zero; the reference subsequently changes to 5 m at approximately 2.5 s. The trace therefore combines startup and a reference change.
\fig{../../data/session2/simulink_scope.png}{.80}{Original supplied scope with reference and altitude traces. The initial output is not settled at the initial reference.}{fig:originalscope}
The regenerated 0 to 5 m step in \Fref{fig:nominal} avoids combining these initial conditions. The included \texttt{matlab/build\_simulink\_model.m} creates a native model with explicit 0--5 V command saturation when executed in Simulink. The original native \texttt{.slx} file was not supplied.
\page

\section{Complete Circuit and Interface}\label{app:circuit}
\subsection{Signal Stages}
\Fref{fig:circuit} uses the component numbering of the supplied circuit files. Its bias input connects to the P1 wiper in \Fref{fig:biasnetwork}; the two figures together define the complete controller signal and bias networks.
\begin{figure}[H]\centering\resizebox{.94\linewidth}{!}{\input{\ReportRoot sections/circuit.tex}}
\caption{Labelled three-stage proportional controller. DAC-0 feeds the inverting sensor branch; DAC-1 feeds the non-inverting reference branch. ADC-2 monitors the sensor and ADC-1 receives the command.}\label{fig:circuit}
\end{figure}
All three op-amps use pin 7 for $+15$ V, pin 4 for $-15$ V, pin 2 for the inverting input, pin 3 for the non-inverting input and pin 6 for output. Pins 1 and 5 (offset null) and pin 8 (NC) are unused in the documented 741 design. \Tref{tab:daq} gives the interface mapping.
\page

\subsection{Bias Network and Supply Decoupling}
\begin{figure}[H]\centering
\begin{circuitikz}[american,scale=.85,transform shape]
\draw (0,0) node[ground]{} to[R,l=$R_{11}\ 3.9\,\mathrm{k}\Omega$] (0,-1.5) coordinate(A)
 to[R,l={$P_{\rm top}=10W\,\mathrm{k}\Omega$}] (0,-3) coordinate(W)
 to[R,l={$P_{\rm bot}=10(1-W)\,\mathrm{k}\Omega$}] (0,-4.5) coordinate(B)
 to[R,l=$R_{13}\ 22\,\mathrm{k}\Omega$] (0,-6) node[below]{$-15$ V};
\draw (A)--(-3,-1.5) to[R,l_=$R_{12}\ 1.2\,\mathrm{k}\Omega$] (-3,-4.5)--(B);
\draw (W) node[circ]{} -- (5,-3) -- (5,-3.7) node[right]{P1 wiper to $R_9$ / U3 pin 2};
\node[left] at (A){A};\node[left] at (B){B};
\node[align=left,anchor=west] at (4.1,-4.9){P1 total: 10 k$\Omega$\\$W=0.783081$ nominal\\wiper loaded through U3};
\draw (7,0) node[above]{$+15$ V (pin 7)} to[C,l=$100\,\mathrm{nF}$] (7,-1.4) node[ground]{};
\draw (10,0) node[above]{$-15$ V (pin 4)} to[C,l=$100\,\mathrm{nF}$] (10,-1.4) node[ground]{};
\node[align=center] at (8.5,-2.7){repeat at U1, U2 and U3\\six capacitors in the design};
\end{circuitikz}
\caption{Bias potentiometer represented by its two segments, shunt resistor and end resistors; supply decoupling is repeated at each op-amp.}\label{fig:biasnetwork}
\end{figure}
\begin{table}[H]\centering\small
\caption{DAQ and supply connections.}\label{tab:daq}
\begin{tabularx}{\linewidth}{lX}\toprule\thead
Connection & Controller node\\\midrule
DAC-1 & $r_v$, U1 reference branch through $R_1$.\\
DAC-0 & $y_v$, U1 sensor branch through $R_3$.\\
ADC-2 & Sensor monitor at $y_v$.\\
ADC-1 & Command $u$, U3 pin 6; permitted plant range 0--5 V.\\
Ground & Shared by DAQ, bias network and bench supply.\\
$\pm15$ V & Op-amp rails; these do not define the permitted plant command.\\\bottomrule
\end{tabularx}\end{table}
The LTspice source and supplied circuit image place $R_{11}$ at ground. The supplied KiCad file labelled that endpoint $+15$ V; its bundled copy is corrected to ground to match those two records. This correction is documented in \texttt{hardware/design\_basis.md}. Calibration sets $r_v=y_v$, then adjusts P1 until the loaded ADC-1 command is 2.500 V. Keep the wiper fixed during testing.
\page

\subsection{Bias Loading and Stage Gain Calculations}\label{app:loading}
With the wiper unloaded, $R_P=P1\parallel R_{12}=10\parallel1.2=1.0714286$ k$\Omega$. The end-to-end current and pot-end voltages are
\begin{align}
I&=\frac{-15}{3.9+1.0714286+22}=-0.556143\unit{mA},\\
V_A&=I(3.9)=-2.168962\unit{V},\qquad V_B=I(3.9+1.0714286)=-2.764831\unit{V},\\
V_{\rm th}&=V_A+W(V_B-V_A).
\end{align}
Set the supply sources to zero and inject a 1 A test current at the wiper. The node-conductance equations for A, B and the wiper are
\begin{equation}
\begin{bmatrix}
G_{11}+G_{12}+G_t&-G_{12}&-G_t\\
-G_{12}&G_{13}+G_{12}+G_b&-G_b\\
-G_t&-G_b&G_t+G_b
\end{bmatrix}
\begin{bmatrix}V_A'\\V_B'\\V_w'\end{bmatrix}
=\begin{bmatrix}0\\0\\1\end{bmatrix},
\end{equation}
where $G_i=1/R_i$, $G_t=1/(WP1)$ and $G_b=1/[(1-W)P1]$. The test voltage $V_w'$ gives $R_{\rm th}$. Solving the hover equation
\begin{equation}
-\frac{R_{10}}{R_9+R_{\rm th}(W)}V_{\rm th}(W)=2.5
\end{equation}
gives $W=0.783080594$, $R_{\rm th}=5423.0014$ $\Omega$ and $V_{\rm th}=-2.6355750$ V. The loaded zero-error wiper is $-2.500$ V.

For U3, define $r=R_{10}/(R_9+R_{\rm th})=0.94855960$. Since the non-inverting divider gives $V_+=V_2/2$, virtual short and KCL give
\begin{align}
u&=(1+r)V_+-rV_{\rm th}=\frac{1+r}{2}V_2-rV_{\rm th},\\
g_3&=\frac{1+0.94855960}{2}=0.97427980,\\
K_p&=\left(\frac{1}{100}\right)\left(1+\frac{4.7}{10}\right)(0.97427980)=0.014321913.
\end{align}
The ideal buffered result 0.0147 differs by 2.57\%, while the selected model value 0.01432 differs from the loaded calculation by only 0.0134\%. These are circuit-model calculations; a measured input/output slope of the final board was not supplied.
\page

\section{Detailed Calculations}\label{app:calculations}
\input{\ReportRoot sections/model_calculations.tex}
\page
\subsection{Identification and Timing}
\Tref{tab:rawsamples} gives selected original samples and derived velocities. The finite manual step is evident in the input column. Endpoint times are listed for consistency with the CSV; each derivative physically belongs to its interval midpoint.
\begin{table}[H]\centering\small
\caption{Selected identification samples and backward-difference results.}\label{tab:rawsamples}
\begin{tabular}{rrrrr}\toprule\thead
$t$ [s] & $u$ [V] & $y_v$ [V] & $y$ [m] & $v_i$ [m/s]\\\midrule
15.7 & 2.495 & 0.000 & 0.000 & 0.00\\
15.8 & 2.908 & 0.000 & 0.000 & 0.00\\
15.9 & 4.607 & 0.000 & 0.000 & 0.00\\
16.0 & 5.000 & 0.034 & 0.037 & 0.37\\
16.9 & 5.000 & 2.098 & 2.256 & 4.08\\
17.0 & 5.000 & 2.513 & 2.703 & 4.47\\
23.9 & 5.000 & 10.000 & 104.823 & 22.36\\
24.0 & 5.000 & 10.000 & 107.076 & 22.53\\
24.1 & 5.000 & 10.000 & 109.345 & 22.69\\\bottomrule
\end{tabular}\end{table}
\begin{align}
v(24.0)&=(107.076-104.823)/0.1=22.53\unit{m/s},\\
v(24.1)&=(109.345-107.076)/0.1=22.69\unit{m/s},\\
\bar v_\infty&=\frac1{404}\sum_{t_i\ge120}\frac{y_i-y_{i-1}}{0.1}=35.68747525\unit{m/s},\\
A&=(35.68747525-0)/2.5=14.27499010\unit{m/(s\,V)},\\
v_{63}&=0.632(35.68747525)=22.55448436\unit{m/s},\\
t_{63}&=24.0+0.1\frac{22.55448436-22.53}{22.69-22.53}=24.01530272\unit{s},\\
\tau_{63}&=24.01530272-15.9=8.11530272\unit{s}.
\end{align}
The exponential fit minimises $\sum_{17\le t_i\le60}[v_i-V_\infty(1-e^{-(t_i-t_0)/\tau})]^2$. Its outputs are $V_\infty=35.68747059$ m/s, $\tau=8.10393362$ s and $t_0=15.91583783$ s. These two methods corroborate rounding to $\tau=8.1$ s. They do not create two experimental trials.

The sensor least-squares fit in \Eref{eq:sensor} excludes zero and clipped pairs. The check $k_s(2.703)=0.93(2.703)=2.51379$ V agrees with the recorded 2.513 V within rounding.
\page

\input{\ReportRoot sections/root_locus_calculations.tex}
\subsection{Root-Locus Angle and Magnitude Checks}
\textbf{Step 6 -- Verify the pole position.}
The locus for variable $K_p$ uses $G_0(s)=k_sG_y(s)=(Ak_s/\tau)/[s(s+1/\tau)]$, where $Ak_s/\tau=1.638959$ s$^{-2}$. At the selected upper pole $s_d=-0.0617284+j0.1402123$,
\begin{align}
\angle s_d&=113.76147^\circ,\qquad
\angle(s_d+1/\tau)=66.23853^\circ,\\
\angle G_0(s_d)&=-113.76147^\circ-66.23853^\circ=-180^\circ,\\
K_p&=\frac1{|G_0(s_d)|}
=\frac{|s_d|\,|s_d+1/\tau|}{Ak_s/\tau}
=\frac{(0.153198845)^2}{1.63895852}=0.01432.
\end{align}
Both the negative-feedback angle criterion and magnitude criterion are satisfied. The quadratic roots independently give the same pole pair.

\page
\subsection{Nominal Step and Settling Calculation}
\textbf{Step 7 -- Evaluate the time response.}
Substituting into \Eref{eq:exactstep},
\begin{equation}
y(t)=5\left[1-e^{-0.061728395t}
\left(\cos(0.140212308t)+0.44025\sin(0.140212308t)\right)\right].
\end{equation}
The first peak occurs at $t_p=\pi/\omega_d=22.405969$ s. Since $\cos\pi=-1$ and $\sin\pi=0$,
\begin{equation}
y(t_p)=5[1+e^{-0.061728395(22.405969)}]=6.254019\unit{m},
\end{equation}
so $M_p=(6.254019-5)/5\times100=25.080375\%$.

\textbf{Step 8 -- Calculate settling.} The error envelope satisfies $|y-5|/5\le e^{-\sigma t}/\sqrt{1-\zeta^2}$. Solving this bound for 2\% gives
\begin{equation}
T_{\rm env}=-\frac{\ln(0.02\sqrt{1-0.4029299^2})}{0.061728395}=64.80975\unit{s}.
\end{equation}
Actual settling is found by evaluating the full response, locating the last sample for which $|y-5|>0.1$ m, and taking the next sample. A 0.005 s grid gives 54.89 s. This definition agrees with reference-based settling calculations \cite{stepinfo}.

\textbf{Step 9 -- Check command headroom.} Command headroom follows from $u=2.5+K_pk_s(5-y)$: the highest command is 2.566588 V at the step, and the lowest is $2.5-0.01432(0.93)(1.254019)=2.483299$ V at the altitude peak. Neither reaches the 0 or 5 V limit.
\page

\subsection{Aerodynamic and Component Uncertainty}
With fixed $k_T$, the uncertain characteristic polynomial is $s^2+b's+k_TK_pk_s$. The worst scalar case has $b'=0.9b_0$ and $K_p'=1.1K_{p,0}$. Applying the required damping ratio gives
\begin{equation}\label{eq:jointbound}
K_{p,0}<\frac{(0.9b_0)^2}{4k_Tk_s\zeta_{\min}^2(1.1)}
=\frac{0.81}{1.1}(0.01815443)=0.01336826.
\end{equation}
At the supplied gain the worst-case damping is $(0.9/\sqrt{1.1})(0.4029299)=0.34575$, and \Eref{eq:overshoot} gives 31.42\% overshoot. This explains why passing drag-only uncertainty is insufficient for the joint case.

For independent loop-resistor tolerances, use the U1 coefficients and U2/U3 gains defined in \Sref{sec:hardware}. Assign each of $R_1$--$R_{10}$ either 0.9 or 1.1 times its nominal value, then use both drag extremes. There are $2^{10}\times2=2048$ cases. The bias network is held at its nominal Thevenin resistance and zero-error output is recalibrated in each case. \Tref{tab:corners} reports this optimistic subset.
\begin{table}[H]\centering\small
\caption{Independent loop-resistor corner predictions.}\label{tab:corners}
\begin{tabularx}{\linewidth}{Xr}\toprule\thead
Quantity & Range / worst value\\\midrule
Feedback gain $K_f$ & 0.009072--0.022784\\
Reference/feedback ratio $a_r/a_f$ & 0.672104--1.487864\\
Absolute DC tracking error & 48.7864\%\\
Overshoot about each case's own final altitude & 38.9461\%\\\bottomrule
\end{tabularx}\end{table}
For example, $R_1=0.9(100\,\mathrm{k}\Omega)$, $R_2=1.1(1\,\mathrm{k}\Omega)$, $R_3=1.1(100\,\mathrm{k}\Omega)$ and $R_4=0.9(1\,\mathrm{k}\Omega)$ give
\begin{equation}
a_f=0.9/110=0.00818182,\quad
\frac{a_r}{a_f}=\frac{(1+0.00818182)(1.1)/(90+1.1)}{0.00818182}=1.487864.
\end{equation}
Thus a 5 m reference tends to 7.43932 m even after hover recalibration. The type number cannot remove unequal reference and feedback scaling.
\page
\subsection{Practical Gain and Ratio Correction}
The corner distributions in \Fref{fig:components} separate dynamic overshoot from DC ratio error.
\fig{component_robustness}{.92}{Loop-resistor corners. Dashed limits are 30\% overshoot and 8\% tracking error. The histograms show corner counts, not probabilities.}{fig:components}
Failure in this subset rejects a guarantee for all components. Bias-network and op-amp variations can add error. After matching the U1 ratios, a practical gain correction is to replace both 1 k$\Omega$ resistors with matched 680 $\Omega$ values:
\begin{align}
K_{p,\rm new}&=\frac{680}{100000}(1.47)(0.97427980)=0.009738901,\\
\zeta_{\rm worst}&=\frac{0.9b_0}{2\sqrt{k_T(1.1K_{p,\rm new})k_s}}\approx0.4193,\\
M_{p,\rm worst}&\approx23.44\%,\\
|\delta u_b|_{\max}&=0.08(5)(0.009738901)(0.93)=3.623\unit{mV}.
\end{align}
This scalar check assumes matched reference and feedback ratios; it is not a guarantee for independent resistor corners. Lower gain increases damping but tightens the hover-bias allowance. The board must be recalibrated and retested after any change.
\page

\subsection{Experimental Metrics and Uncertainty}
The digitised plot uses pixel bounds $x=602\ldots1414$ for 0--120 s and $p_y=60\ldots540$ for 20--0 m. Therefore
\begin{align}
t&=\frac{x-602}{1414-602}(120),\qquad h=\frac{540-p_y}{540-60}(20),\\
\Delta t&=0.147783\unit{s/pixel},\qquad\Delta h=0.041667\unit{m/pixel}.
\end{align}
The extracted peak is 6.229167 m at 22.019704 s. With the displayed reference, terminal height and zero starting altitude,
\begin{align}
e_f&=5.033-4.897=0.136\unit{m},\\
\eta_{\rm track}&=100(1-0.136/5.033)=97.297834\%,\\
M_{p,\rm ref}&=100(6.229167-5.033)/5.033=23.766475\%,\\
\delta M_{p,\rm ref}&\approx100(0.041667)/5.033=0.82787\text{ percentage points}.
\end{align}
This practical overshoot metric uses the commanded reference. If the terminal display is used as a proxy for the actual settled output, the conventional final-output-normalised overshoot is $(6.229167-4.897)/4.897\times100\approx27.20\%$. Both estimates are below 30\%, but the screenshot alone does not establish an exact asymptotic output.
\Tref{tab:comparison} compares the analytical, simulated and experimental metrics.
\begin{table}[H]\centering\small
\caption{Comparison of calculation, simulation and retained experiment.}\label{tab:comparison}
\begin{tabularx}{\linewidth}{lXXX}\toprule\thead
Metric & Analytical & Simulation & Experiment\\\midrule
Reference & 5.000 m & 5.000 m & 5.033 m\\
Peak & 6.25402 m & 6.25402 m & $\approx6.229$ m\\
Peak time & 22.40597 s & 22.405 s & $\approx22.020$ s\\
Reference overshoot & 25.0804\% & 25.0804\% & $\approx23.7665\%$\\
Tracking accuracy & 100\% ideal & $\to100\%$ & 97.2978\% displayed\\
2\% reference settling & 64.81 s envelope & 54.89 s & Not established\\\bottomrule
\end{tabularx}\end{table}
For the 5.033 m reference the 8\% band is $5.033(1\pm0.08)=4.63036$--5.43564 m. The 2\% band is 4.93234--5.13366 m. The terminal height lies inside the first but outside the second. The discrepancy cannot be removed by redefining the settling band around the biased terminal value while claiming reference tracking.

The command-bias calculation uses \Eref{eq:bias}: $\delta y=0.002/(0.01432\times0.93)=0.150177$ m. The illustrative test in \Fref{fig:disturbance} injects a positive 2 mV bias at 70 s; it is a model prediction rather than a recorded disturbance experiment.
\page

\section{Additional Validation Plots}\label{app:validation}
\Fref{fig:posval} checks the integrated position response; \Fref{fig:aero} compares the three aerodynamic cases listed in \Tref{tab:aero}. The underlying values are retained as CSV and JSON in the \texttt{results} folder.
\fig{position_validation}{.88}{Position-domain model comparison and residuals for the retained identification run.}{fig:posval}
\fig{aerodynamic_robustness}{.88}{Predicted altitude responses for $b=0.9b_0,b_0,1.1b_0$, with local thrust slope held fixed.}{fig:aero}
\page
\Fref{fig:disturbance} demonstrates sensitivity to a defined input bias. At nominal gain its predicted final displacement is +0.150 m, within the 0.4 m allowance for a 5 m reference.
\fig{disturbance_prediction}{.95}{Illustrative simulated response to +2 mV actuator bias added at 70 s. The shaded region is the 8\% tracking band.}{fig:disturbance}
This simulation tests one stated disturbance amplitude. It does not imply rejection of arbitrary disturbance magnitudes or establish experimental disturbance performance.
\page

\section{Identification Data and Code}\label{app:data}
\subsection{Data Provenance}
The unmodified \texttt{data/session1/PTRJUN003\_raw.csv} contains time, command voltage, sensor voltage, altitude and stored speed. The processed copy excludes the 31 pre-reset rows and adds recalculated velocity, midpoint time and model predictions.

\Tref{tab:rawsamples} gives selected calculations. The complete CSV and executable analysis reproduce all estimates; no previous-year data replace missing trials.


\subsection{Organised Source Code}\label{app:code}
The numerical pipeline is separated by task. All stages read the shared \RepoFile{config/lab_parameters.json}{parameter file}; output CSVs and metrics are stored in \RepoLink{results}{results}. Python generated the report figures. MATLAB equivalents require MATLAB with Control System Toolbox; the Simulink builder additionally requires Simulink.
\begin{table}[H]\centering\small
\caption{Code modules and their calculation responsibilities.}\label{tab:modules}
\begin{tabularx}{\linewidth}{p{43mm}X}\toprule\thead
Source location & Responsibility\\\midrule
\RepoLink{analysis/modules}{Python modules} & Identification, controller response/root locus, robustness, hardware digitisation and disturbance simulation.\\
\RepoLink{matlab/identification}{MATLAB identification} & Backward differences, velocity gain, 63.2\% timing and sensor calibration.\\
\RepoLink{matlab/design}{MATLAB design} & Root-locus geometry, angle/magnitude checks, gain limits and response plots.\\
\RepoLink{matlab/validation}{MATLAB validation} & Independent nominal and aerodynamic controller checks.\\
\RepoLink{matlab/simulink}{Simulink builder} & Reconstructs and saves the nominal model.\\
\RepoLink{calculations}{Worked calculations} & Separate calculation PDF, LaTeX equations and topic index.\\\bottomrule
\end{tabularx}\end{table}
Run the Python launcher from the project root:
\lstset{basicstyle=\ttfamily\small,breaklines=true,frame=single,rulecolor=\color{AngelBlue},columns=fullflexible,showstringspaces=false}
\begin{lstlisting}[language=bash]
python -m pip install -r requirements.txt
python analysis/analyse.py
\end{lstlisting}
In MATLAB, set the Current Folder to \texttt{matlab}, then enter:
\begin{lstlisting}[language=Matlab]
run_all
% Or run only the detailed root locus method:
root_locus_design
\end{lstlisting}
\RepoFile{CODE_GUIDE.md}{The code guide} records dependencies, file responsibilities and step conventions. Voltage identification uses a 2.5 to 5.0 V step ($\Delta u=2.5$ V); the nominal controller test uses a 0 to 5 m altitude reference.
\page
\section{Repository and Reproduction}\label{app:repository}
\RepoStatus\ Open the \href{\RepoRoot}{complete GitHub repository} or use the folder links in \Tref{tab:repo}. Figure, table, equation and appendix references remain internal PDF hyperlinks. The repository guide and calculation index explain how to reproduce each result.
\begin{table}[H]\centering\small
\caption{Clickable evidence map and corresponding Git folders.}\label{tab:repo}
\begin{tabularx}{\linewidth}{p{37mm}X}\toprule\thead
Folder / evidence link & Contents\\\midrule
\RepoLink{data/session1}{\texttt{data/session1}} & Original identification CSV; source description in \Aref{app:data}.\\
\RepoLink{data/session2}{\texttt{data/session2}} & Original experimental display and Simulink screenshots.\\
\RepoLink{hardware}{\texttt{hardware}} & Board photograph, LTspice/KiCad design sources and circuit notes.\\
\RepoLink{analysis}{\texttt{analysis}} & Python identification, simulation, digitisation and tolerance calculations.\\
\RepoLink{matlab}{\texttt{matlab}} & Controller verification and Simulink model builder.\\
\RepoLink{results}{\texttt{results}} & Processed CSVs, nominal traces, corners and scalar metrics.\\
\RepoLink{report}{\texttt{report}} & LaTeX source, figures and compiled report.\\
\RepoLink{scripts}{\texttt{scripts}} & Build, link configuration and folder indexes.\\
\RepoLink{config}{\texttt{config}} & One shared parameter file for MATLAB and Python.\\
\RepoLink{calculations}{\texttt{calculations}} & Worked calculation PDF and equation source.\\\bottomrule
\end{tabularx}\end{table}

\subsection{Build the Report}
Extract the ZIP, keep its folder structure, install the Python requirements, and run:
\begin{lstlisting}[language=bash]
python -m pip install -r requirements.txt
bash scripts/build.sh
\end{lstlisting}
The build regenerates numerical results and figures, then compiles \texttt{report/main.tex} with \texttt{latexmk -pdf}. A LaTeX distribution with TikZ and Circuitikz is required. The package also contains the compiled PDF for immediate viewing.

\page
\subsection{GitHub Folder Links}
The report uses the repository supplied for this project:
\begin{center}\href{\RepoRoot}{NateDogg247 / EEE3094S---Lab-2-PTRJUN003-}\end{center}
The shared branch setting is \texttt{main}. Supporting-folder URLs use \texttt{/tree/main/}; individual source-file links use \texttt{/blob/main/}. \texttt{report/repository.tex} stores these settings centrally. Recompile the report after changing that file.

\subsection{Page Count and Reproduction}
The main report runs from Introduction on page 1 to Final Remarks on page 14. The title, table of contents, lists of figures and tables, references and appendices are excluded from this count. References use R-prefixed page numbers and appendices use A-prefixed page numbers.

The numerical pipeline regenerates identification, root-locus, response, tolerance and experimental metrics. The report and calculation supplement retain the corresponding worked equations. The available record contains one identification run and one controller display; it cannot establish repeated-trial averages or an assessor-awarded demonstration mark.
